\subsection{Groverio algoritmo idėja}
\label{sec:groverio-algoritmo-ideja}

Groverio algoritmas sprendžia nestruktūrizuotos paieškos problemą: duotoje $N$ elementų erdvėje reikia rasti pažymėtą elementą. Klasikinėje paieškoje blogiausiu atveju reikia patikrinti $O(N)$ elementų, o Groverio algoritmas pažymėtą elementą randa naudojant apie $O(\sqrt{N})$ oracle užklausų \cite{grover-1996,nielsen-chuang-2010}.

Algoritmas prasideda nuo vienodos superpozicijos paruošimo. Jeigu paieškos registre yra $n$ kubitų, po Hadamard vartų visoms paieškos kubitų pozicijoms gaunama būsena:

\begin{equation}
    \frac{1}{\sqrt{N}} \sum_{x=0}^{N-1} \lvert x \rangle,
    \quad N = 2^n.
\end{equation}

Tada kartojami du pagrindiniai veiksmai:

\begin{itemize}
    \item \textbf{oracle} -- operacija, kuri pažymi ieškomą būseną. Dažnai tai aprašoma kaip pažymėtos būsenos fazės pakeitimas, t. y. ieškomos būsenos amplitudė dauginama iš $-1$;
    \item \textbf{difuzijos operatorius} -- operacija, vadinama amplitudžių inversija apie vidurkį. Ji padidina pažymėtos būsenos amplitudę ir sumažina kitų būsenų amplitudes.
\end{itemize}

Vienas oracle ir difuzijos operatoriaus pritaikymas sudaro vieną Groverio iteraciją. Kai pažymėtas vienas elementas, apytikslis iteracijų skaičius yra:

\begin{equation}
    k \approx \frac{\pi}{4}\sqrt{N}.
\end{equation}

Po iteracijų registras matuojamas arba analizuojamos jo amplitudės. Idealiu atveju pažymėtos būsenos tikimybė tampa didelė, todėl matavimas dažniausiai grąžina ieškomą elementą. Ši savybė daro Groverio algoritmą tinkamą šiam darbui: jis pakankamai mažas, kad būtų realizuojamas kaip vartotojo erdvės programa, bet kartu naudoja kelis svarbius kvantinius veiksmus -- superpoziciją, oracle, daugkartinę iteraciją ir tikimybių analizę.
